\subsection{Premeasures and measures}

Let T be a topological space, and let \(\mathcal{K}\) be the set of compact subsets of T, ordered by inclusion. For every \(K\in\mathcal{K}\) , let \(\mathcal{M}(K;C)\) be the set of complex measures on K. For every pair (K,L) of elements of \(\mathcal{K}\) such that \(K\subset L\) , let \(\iota_{KL}\) be the mapping of \(\mathcal{M}(L;C)\) into \(\mathcal{M}(K;C)\) that associates to each measure \(\mu\) on L the measure \(\mu_{K}\) induced by \(\mu\) on K (Ch. IV, §5, No.~7, Def. 4). Then \(\iota_{KM}=\iota_{KL}\circ\iota_{LM}\) when K, L and M are compact subsets of T such that \(K\subset L\subset M\) ; this follows from the transitivity of induced measures (Ch. V, §7, No.~2, Prop. 4). The elements of the inverse limit of the family \((\mathcal{M}(K;C))_{K\in\mathcal{K}}\) for the mappings \(\iota_{KL}\) will be called premeasures on T. In other words:

DEFINITION 3. --- One calls premeasure on a topological space T every mapping w that associates, to every compact subset K of T, a measure \(w_{K}\) on K, and that has the following property:

If K and L are compact subsets of T such that \(K \subset L\) , the measure \((w_{\mathrm{L}})_{\mathrm{K}}\) induced by \(w_{L}\) on K is equal to \(w_{K}\) .

The premeasure w is said to be real (resp. positive) if all of the measures \(w_{K}\) are real (resp. positive).

Let \(w\) and \(w'\) be two premeasures on \(T\) , \(t\) a complex number; the pre-measures \(w + w'\) and \(tw\) are defined by the formulas \((w + w')_K = w_K + w'_K\) , \((tw)_K = tw_K\) for every compact subset \(K\) of \(T\) . The premeasures on \(T\) obviously form a vector space, which is denoted \(\mathcal{P}(T; \mathbf{C})\) ; the space of real premeasures will be denoted \(\mathcal{P}(T; \mathbf{R})\) , or more often \(\mathcal{P}(T)\) , and the convex cone of positive premeasures will be denoted \(\mathcal{P}_+(T)\) . Let \(w\) be a premeasure; the mapping \(K \mapsto |w_K|\) is then a premeasure on \(T\) (Ch. IV, §5, No.~7, Lemma 3), which will be denoted \(|w|\) . If \(w\) is real, one sets \(w^{+} = \frac{1}{2} (|w| + w)\) , \(w^{-} = \frac{1}{2} (|w| - w)\) ; these two premeasures being positive, one sees that every real premeasure is the difference of two positive premeasures. Clearly \((w^{+})_{\mathrm{K}} = (w_{\mathrm{K}})^{+}\) , \((w^{-})_{\mathrm{K}} = (w_{\mathrm{K}})^{-}\) for every compact subset K of T.

The vector space \(\mathcal{P}(\mathrm{T})\) is ordered by the cone \(\mathcal{P}_{+}(\mathrm{T})\) . It is clear that \(w^{+} = \sup(w,0)\) , \(w^{-} = \sup(-w,0)\) ; consequently, \(\mathcal{P}(\mathrm{T})\) is lattice-ordered and \(\sup(w,w') = w + (w' - w)^{+}\) , \(\inf(w,w') = w - (w' - w)^{-}\) . Moreover, clearly

\[
\big (\sup (w, w ^ {\prime}) \big) _ {\mathbf {K}} = \sup (w _ {\mathbf {K}}, w _ {\mathbf {K}} ^ {\prime}), \quad \big (\inf (w, w ^ {\prime}) \big) _ {\mathbf {K}} = \inf (w _ {\mathbf {K}}, w _ {\mathbf {K}} ^ {\prime})
\]

for every compact subset K of T.

DEFINITION 4. --- Let w be a positive premeasure on T. We shall set, for every function \(f \in \mathcal{F}_{+}(T)\) ,

\[
w ^ {\bullet} (f) = \sup _ {K} \left(w _ {K}\right) ^ {\bullet} \left(f _ {K}\right),\tag{1}
\]

where K runs over the set of compact subsets of T.

For each compact set K, let \(p^{K}\) be the image encumbrance of the encumbrance \((w_{\mathrm{K}})^{\bullet}\) under the canonical injection of K into T; \(w^{\bullet}\) is the upper envelope of the encumbrances \(p^{K}\) , hence is an encumbrance (No.~1, Prop. 1). One says that \(w^{\bullet}\) is the essential upper integral associated with the positive premeasure w. One often writes \(\int^{\bullet} f dw\) or \(\int^{\bullet} f(t) dw(t)\) instead of \(w^{\bullet}(f)\) .

Remark 1). --- If \(v\) and \(w\) are two positive premeasures, then \((v + w)^{\bullet} = v^{\bullet} + w^{\bullet}\) (Ch. V, §1, No.~1, Prop. 3). If \(v\) and \(w\) are two complex premeasures, then \(|v + w|^{\bullet} \leqslant |v|^{\bullet} + |w|^{\bullet}\) .

PROPOSITION 2. --- a) Let w be a positive premeasure. For every compact subset K of T, the encumbrance \((w^{\bullet})_{\mathrm{K}}\) induced by \(w^{\bullet}\) on K is equal to \((w_{\mathrm{K}})^{\bullet}\) . For every function \(f \in \mathcal{F}_{+}(\mathrm{T})\) , one has the relations \((w_{\mathrm{K}})^{\bullet}(f_{\mathrm{K}}) = w^{\bullet}(f\varphi_{\mathrm{K}})\) and

\[
w ^ {\bullet} (f) = \sup _ {K} w ^ {\bullet} (f \varphi_ {K}).\tag{2}
\]

\begin{enumerate}
\def\labelenumi{\alph{enumi})}
\setcounter{enumi}{1}
\tightlist
\item
  Conversely, let \(p\) be an encumbrance on \(T\) satisfying the following conditions:
\end{enumerate}

\begin{enumerate}
\def\labelenumi{\arabic{enumi})}
\item
  For every compact subset K of T, there exists a positive measure \(w_{\mathbf{K}}\) on K such that \(p_{\mathbf{K}} = (w_{\mathbf{K}})^{\bullet}\) .
\item
  For every function \(f \in \mathcal{F}_{+}(\mathrm{T})\) , \(p(f) = \sup_{\mathrm{K}} p(f\varphi_{\mathrm{K}})\) .
\end{enumerate}

The mapping \(w: \mathbf{K} \mapsto w_{\mathbf{K}}\) is then a positive premeasure on \(\mathbf{T}\) , and \(p = w^{\bullet}\) .

Let us prove \(a\) : let \(g \in \mathcal{F}_{+}(K)\) and let \(g^{0}\) be the extension by zero of \(g\) to \(T\) ; then, by the definition of induced encumbrances,

\[
(w ^ {\bullet}) _ {K} (g) = w ^ {\bullet} \left(g ^ {0}\right) = \sup _ {L} \left(w _ {L}\right) ^ {\bullet} \left(g ^ {0} \mid L\right),
\]

where L runs over the set of compact subsets of T, or merely the set of those that contain K. But if L contains K, then \((w_{\mathrm{L}})^{\bullet}(g^{0}|\mathrm{L})=(w_{\mathrm{K}})^{\bullet}(g)\) from the fact that \(g^{0}|L\) is zero outside K (Ch. V, §7, No.~1, Prop. 1), which proves the first assertion. Therefore

\[
(w _ {K}) ^ {\bullet} (f _ {K}) = (w ^ {\bullet}) _ {K} (f _ {K}) = w ^ {\bullet} \left((f _ {K}) ^ {0}\right) = w ^ {\bullet} (f \varphi_ {K})
\]

for all \(f \in \mathcal{F}_{+}(\mathrm{T})\) , and (2) merely translates the formula (1).

Let us pass to \(b\) : the measure \(w_{\mathbf{K}}\) considered in 1) is unique (Ch. V, §1, No.~1). Let us show that the mapping \(\mathbf{K} \mapsto w_{\mathbf{K}}\) is a premeasure: let \(\mathbf{K}\) and \(\mathbf{L}\) be two compact subsets such that \(\mathbf{K} \subset \mathbf{L}\) , and let \(\lambda\) be the measure induced by \(w_{\mathbf{L}}\) on \(\mathbf{K}\) ; everything comes down to showing that \(\lambda^{\bullet} = (w_{\mathbf{K}})^{\bullet}\) . Now, \(\lambda^{\bullet} = ((w_{\mathbf{L}})^{\bullet})_{\mathbf{K}}\) (Ch. V, §7, No.~1, Prop. 1); since \((w_{\mathbf{L}})^{\bullet} = p_{\mathbf{L}}\) , we have \(\lambda^{\bullet} = (p_{\mathbf{L}})_{\mathbf{K}} = p_{\mathbf{K}} = (w_{\mathbf{K}})^{\bullet}\) .

Let us denote by \(w\) the premeasure \(K \mapsto w_K\) ; since \(p_K = (w_K)^{\bullet} = (w^{\bullet})_K\) , we have \(p(f\varphi_K) = p_K(f_K) = (w^{\bullet})_K(f_K) = w^{\bullet}(f\varphi_K)\) . The two encumbrances \(p\) and \(w^{\bullet}\) are therefore equal by virtue of the formula (2) and the hypothesis 2) on \(p\) .

Q.E.D.

Since the induced encumbrance \((w^{\bullet})_{\mathbf{K}}\) is equal to \((w_{\mathbf{K}})^{\bullet}\) , there is no ambiguity in simply writing \(w_{\mathbf{K}}^{\bullet}\) . We shall employ this notation from now on.

COROLLARY. --- Let v and w be two positive premeasures on T, such that \(v^{\bullet}(\mathbf{L}) = w^{\bullet}(\mathbf{L})\) for every compact subset L of T; then v = w. In particular, the relation \(v^{\bullet} = w^{\bullet}\) implies v = w.

For, let K be a compact set in T; for every compact set \(L \subset K\) , one has the relation

\[
w _ {\mathrm{K}} (\mathrm{L}) = w _ {\mathrm{K}} ^ {\bullet} (\mathrm{L}) = w ^ {\bullet} (\mathrm{L}) = v ^ {\bullet} (\mathrm{L}) = v _ {\mathrm{K}} ^ {\bullet} (\mathrm{L}) = v _ {\mathrm{K}} (\mathrm{L})
\]

by Prop. 2; therefore \(w_{\mathrm{K}} = v_{\mathrm{K}}\) (Ch. IV, §4, No.~10, Cor. 3 of Prop. 19), and finally \(w = v\) by the definition of premeasures.

DEFINITION 5. --- Let w be a premeasure on a topological space T. One says that w is a measure (resp. a bounded measure) if the encumbrance \(|w|^{\bullet}\) is locally bounded (resp. bounded) (cf.~No.~1, Def. 2).

The set of complex measures on \(\mathbf{T}\) is obviously a vector space (Remark 1), which will be denoted \(\mathcal{M}(\mathbf{T};\mathbf{C})\) . The space of real measures will be denoted \(\mathcal{M}(\mathrm{T};\mathbf{R})\) or more often \(\mathcal{M}(\mathrm{T})\) , and the cone of positive measures will be denoted \(\mathcal{M}_{+}(\mathrm{T})\) .

If w is a complex measure, its real part and its imaginary part are real measures. If w is a real measure, \(w^{+}\) and \(w^{-}\) are positive measures. Every complex (resp. real) measure is thus a linear combination (resp. difference) of positive measures.

Remarks. --- 2) If T is locally compact, then every premeasure w on T is a measure. For, every \(x \in T\) admits a compact neighborhood K, and \(|w|^{\bullet}(K) = \|w_{K}\| < +\infty\) , so that the encumbrance \(|w|^{\bullet}\) is locally bounded.

\begin{enumerate}
\def\labelenumi{\arabic{enumi})}
\setcounter{enumi}{2}
\tightlist
\item
  For every Borel subset A of T (in particular for A = T) and every positive measure \(\mu\) on T, the number \(\mu^{\bullet}(\mathrm{A})\) is the supremum of the measures \(\mu^{\bullet}(\mathrm{K})\) of the compact subsets of A. Indeed, for every compact subset K of A, one has \(\mu^{\bullet}(\mathrm{K}) \leqslant \mu^{\bullet}(\mathrm{A})\) ; on the other hand, if \(\mathfrak{K}\) is the set of compact subsets of T, then
\end{enumerate}

\[
\mu^{\bullet}(\mathrm{A}) = \sup_{\mathrm{K}\in \mathfrak{K}}\mu_{\mathrm{K}}^{\bullet}(\mathrm{A}\cap \mathrm{K}) = \sup_{\mathrm{K}\in \mathfrak{K}}\sup_{\substack{\mathrm{L}\in \mathfrak{K}\\ \mathrm{L}\subset \mathrm{A}\cap \mathrm{K}}}\mu_{\mathrm{K}}^{\bullet}(\mathrm{L})\leqslant \sup_{\substack{\mathrm{L}\in \mathfrak{K}\\ \mathrm{L}\subset \mathrm{A}}}\mu^{\bullet}(\mathrm{L})
\]

by Cor. 1 of Th. 4 of Ch. IV, §4, No.~6.

